\newpage
\recipe{French Onion Soup}{Serves 4}{}{}

\begin{multicols}{2}
	\subsection*{Ingredients}
	\subsubsection*{Soup}
	\begin{ingredients}
		\item 6 yellow onions\footnote{The European variety, which rarely exceeds 3 inches in \diameter , versus the softball+ size of PNW onions. Plan accordingly.}
		\item 55 g butter
		\item 125 ml red wine
		\item 30 ml Cognac (Brandy will work, too)
		\item \sfrac{1}{2} L chicken stock
		\item \sfrac{1}{2} L beef stock
		\item 1 tbsp flour
		\item 1 pinch ground nutmeg
	\end{ingredients}
	
	\columnbreak
	
	\subsection*{}
	\subsubsection*{Topping}
	\begin{ingredients}
		\item 200 g grated Gruyère
		\item Toasted baguette slices
		\item 1 garlic clove
	\end{ingredients}
	\vfill
\end{multicols}

\medskip

% --- Preparation ---
Cut the \textbf{onions} into quarters, then thinly slice them into short strips.
In a large saucepan over medium heat, soften the onions in the \textbf{butter} for 15 minutes.
Continue cooking over high heat for another 15 minutes, or until the onions are well caramelized, stirring constantly and scraping the bottom of the pan.
Deglaze with the \textbf{red wine} and the \textbf{Cognac} (or Brandy), then reduce until almost dry.
Add the \textbf{chicken stock}, \textbf{beef stock}, \textbf{flour}, and \textbf{nutmeg}.
Simmer over medium-high heat for 30 minutes, or until the soup has reduced by about half.
Season with salt and pepper. \\

Ladle the soup into bowls, top each with the garlic-rubbed toasted baguette slices, generously cover with the grated Gruyère, and shortly broil in the oven until the cheese is melted, bubbling, and golden.

\vspace*{0.75 cm}

% --- Description ---
\begin{recipeblock}
	
\end{recipeblock}